\section{Automatizálás és ütemezés}

Ebben a részben a cron és at használatát, valamint a crontab kifejezések magyarázatát találod.

\subsection{Cron alapok}
Cron sor formátum: \texttt{min hour day month weekday command}
\begin{lstlisting}
# Example: every day at 02:00
0 2 * * * /usr/local/bin/backup.sh
# Example: every 15 minutes
*/15 * * * * /usr/local/bin/check.sh
\end{lstlisting}

Használj cron-kalkulátort ellenőrzéshez: \url{https://crontab.guru/}

\subsection{at egyszeri feladat}
\begin{lstlisting}
# Egy egyszeri parancs ma 22:30-kor
echo "/usr/local/bin/job.sh" | at 22:30
# List jobs
atq
# Delete job
atrm 1
\end{lstlisting}

Biztonság:
\begin{itemize}
  \item Kerüljük a jelszavakat és titkos kulcsokat cron bejegyzésekben; használjunk kulcsfájlokat és megfelelő jogosultságokat.
  \item Monitorozzuk a cron futásokat és outputokat (redirect stdout/stderr fájlba) és riasztást építsünk hibákra.
\end{itemize}
